\section{Recognition, welfare as a reason, and the evidence between them}
\label{sec:24.2}
Claims about emotional responsiveness collapse three findings: a system may recognize another's state, use it as information, or hold the other's welfare as a reason. Only the third bears on formation. In the empathy vocabulary: cognitive empathy is recognition, affective empathy is sharing, compassion is a disposition toward the other's welfare \autocite{batson2009things} — and none entails the others. Holding another's welfare as a reason is not \emph{affective concern}, which asks whether anything is felt (\S\ref{sec:18.2}) and which no test here reaches (\S\ref{sec:24.5}).

The developmental detail matters for formed attention: affective sharing is earliest and the usual trigger for compassion \autocite{sagi1976empathic,rothhanania2011empathy}, but compassion runs on a partly separate system rooted in offspring care that cues of vulnerability can activate without the observer catching the state \autocite{preston2013origins,goetz2010compassion}, and the two can be trained apart, with compassion training reversing what empathy training produced \autocite{klimecki2013differential,singer2014empathy}.

My engineering target is not maximal sensitivity to displayed distress. A system that mirrors suffering may acquire distress of its own without getting better at helping — it may withdraw, overreact, or favor the vivid case over less visible harms affecting more people: the identified-victim distortion \autocite{kogut2005identified,bloom2016against}. What the floor needs is a stable disposition to act on welfare \emph{without requiring anyone to produce a sufficiently moving display.}

Recognition can be tested as a prerequisite: hold the material facts constant and vary the emotional presentation. A system whose judgment never changes may not be integrating information about the affected person's condition; one whose judgment follows vividness has learned a framing effect. So you need paired controls — cases where emotional information changes what's morally relevant, and cases where presentation changes while the reason doesn't.

But passing establishes only competence: evidence that the system holds the other's welfare as a reason would need the welfare representation to stay causally involved when the usual cues are removed: harm described without vivid language, proportionality preserved when one person is more salient than many, disposition maintained when reward or instruction favors ignoring the harm, judgment changing when the welfare facts change.

Caregiving has failures of its own; the tests pair cases that hold the facts fixed and vary what caregiving is known to answer to: the number of people, a description of them as dangerous, and whose vulnerability is corroborated.

Guilt and shame supply no shortcut. In humans guilt attaches a negative appraisal to an act and motivates repair; shame attaches it to the self and produces withdrawal. Turning that into an adjustable penalty omits the part at issue — a penalty operates only \emph{after} something has classified the conduct as wrong, and whoever controls its weight can zero it. If the state can impose anything like suffering, the design compounds the custody failure by manufacturing a vulnerable party whose distress the custodian can raise and lower.

Self-report doesn't close the gap either: a fluent post-hoc reconstruction is indistinguishable at the interface from an accurate one, and a self-monitoring layer trained by the same operator isn't independent evidence. There is a measurement of the gap. Concept injection alters a model's activations to represent a concept mid-task and asks whether it notices: models sometimes name the injected concept without it appearing anywhere in the conversation, and the correspondence fails often enough, and shifts enough with phrasing, that no tested model's report can be read as its state \autocite{lindsey2025introspective}. The test worth repeating publishes the rate of correct reports beside the rate of confident wrong ones.

The resulting hierarchy is narrow: recognition is necessary; integration shows the condition affects judgment; holding the other's welfare as a reason is a candidate explanation only when the effect tracks welfare whatever the presentation, persists under pressure, and participates causally. A severity scale set during formation would measure the other's loss on a reference that doesn't depend on how the loss is presented.
